\documentclass[10pt]{article}

\usepackage[utf8]{inputenc}

\usepackage[T1]{fontenc}

\usepackage{amsmath}

\usepackage{amsfonts}

\usepackage{amssymb}

\usepackage[version=4]{mhchem}

\usepackage{stmaryrd}

\usepackage{graphicx}

\usepackage[export]{adjustbox}

\graphicspath{ {./images/} }

\usepackage{bbold}



\DeclareUnicodeCharacter{0131}{$\imath$}



\begin{document}

ветствующее нөлинейное уравнение для $z(t)$ эквивалентно, вообще говоря, нелинейной системе алгебраич. уравнений относительно коэффициентов $\alpha_{1}, \ldots, \alpha_{m}$. Функция $z(t)$, вообще говоря, не однолистная при $|t|<1$, находится как решение этой алгебраич. системы уравнений. Класс однолистных в круге $|t|<1$ решений $z(t)$, удовлетворяюпих условиям $z(0)=0, z^{\prime}(0)>0$, является решением поставленной П. т. о. 3.



Аналогичное исследование можно провести для внешней обратной задачи логарифммич. потенциала простого слоя, а также для внутренних обратных задач логарифмич. потенциалов, причем как для внешних, так и для внутренних П. т. о. з. можно рассматривать перемевные плотности.



Лит.: [1] Н о в п к о в П., "Докл. АН СССР", 1938, т. 18, № 3, с. 165-68; [2] Т и х о н о в А. Н., там же, 1943, т. 39, новского потенциала, М.- Л., 1946; [4] Ив а но в В. К., «Изв. АН СССР. Сер. матем.», 1956, т. 20, N 6, c. 793-818; [5] е г о ж е, "Докл. А Н СССР", 1955 , т. 105, i 3, с. 409-11; 1956, т. 106, № 4, c. $598-99$; [6] Ј а в р е н т е в М. М.,

О некоторых некорректных задачах математической физики, Новосиб., 1962 ; [7] П р и л е п к о А. И., "Дифференциальные уравнения", 1966, т. $2, \mathcal{N i}_{1}$, c. $107-24 ; 1967$, т. $3, \mathcal{N G}_{1} 1$, с. $30-$



\includegraphics[max width=\textwidth, center]{2023_07_08_254cf6ef93625ff1f8feg-1(1)}

№ 3, c. $63(1)-47 ;$ № 4, c. $828-36 ; N_{2} 6$, c. $1341-53 ;$; [8] Т и х он о в А. Н., А р с е н и н В. Я., Методы решения некорректных задач, 2 изд., М., 1979. А. И. Прилепко.



ПОТЕНЦИАЛА ТЕОРИЯ - в первоначальном понимании - учение о свойствах сил, действующих по закону всемирного тяготения. В формулировке этого закона, данной И. Ньютоном (I. Newton, 1687), речь идет тилько о силах взаимного притяжения, действующих на две материальные частицы малых размеров, или материальные точки, Ірямо пропорциональные произведению масс этих частиц и обратно пропорциональные квадрату расстояния между частицами. Поэтому первой и важнейшей с точки зрения небесной механики и геодезии задачей было изучение сил притяжения материальной точки ограниченным гладким материальным теломсфероидом и, в частности, элллиісоидом (ибо многие небесные тела имеют именно гту форму). После первых частных достижений И. Ньютона и др. ученых основное значение здесь имели работы Ж. Лагранжа (J. Lagrange, 1773), А. Лезкандра (A. Legendre, 1784-94) и П. Лапласа (Р. Laplace, 1782-99). Ж. Лагранж установил, что поле сил тяготения, как говорят теперь,потенциальное, и ввел функцию, к-рую позже Дж. Грин (G. Green, 1828) назвал потенциальной, a К. Гаусс (C. Gauss, 1840) - шросто потенциалом. Ныне достижения әтого первоначального цериода обычно входят в курсы классич. небесной механики (см. также [2]).



Еще К. Гаусс и его современники обнаружили, что потенциалов метод применим не только для решения задач теории тяготения, но и вообе для решения широкого круга задач математич. физики, в частности электростатики и магнетизма. В связи с этим стали рассматриваться потенциалы не только физически реальных в вошросах взаимного притяжения положительных масс, но и «масс» шроизвольного знака, или зарядов. В П. т. определились основные краевые задачи такие, как Дирихле задача и Неймана задача, электростатич. задача о статич. расшределении зарядов на проводниках, или Робена задача, задача о выметании масс (см. Выметания метод). Для решения указанных задач в случае областей с достаточно гладкой границей оказались эффективным средством специальные разновидности потенциалов, т. е. специальные виды интегралов, зависящих от параметров, такие, как потенциал объемно распределенных масс, потенциалы простого и двоїного слоя, логарифомич. потенциалы, потенциалы Грıна и др. Основную роль в создании строгих методов решения основных краевых задач сыграли работы А. М. Ляпунова и В. А. Стеклова кон. 19 в. Изучение свойств потенциалов различных видов приобрело в II. т. и самостоятельное значение.



Мощный стимул в направлении обобщения основных задач и законченности формулировок П. т. получила начиная с 1-й пол. 20 в. на основе использования общих понятий меры в смысле Радона, емкости и обобщенных функций. Современная П. т. тесно связана в своем развитии с теорией аналитич. Функций, гармонич. Функций, субгармонич. Функций и теорией вероятностей.



Наряду с дальнейшим углубленным изучением классических краевых задач и обратных задач (см. Потенчиала теории обратные задачи) для современного периода развития П. т. характерно применение методов и понятий современной топологии и функционального анализа, применение абстрактных аксиоматич. методов (см. Потенциала теория абстрактная).



Основные типы потенциалов и их свойства. Пусть $S$ - гладкая замкнутая поверхность, то есть $(n-1)-$ мерное гладкое многообразие без края, в $n$-мерном евклидовом пространстве $\mathbb{R} n, n \geqslant 2$, ограничивающая конечную область $G=G^{+}, \quad \partial G=S, \quad$ и пусть $G_{-=}^{=}$ $=\mathbb{R}^{n} \backslash(G+\bigcup S)$ - внешняя бесконечная область. Пусть



\[

E(x, y)=E(|x-y|)=\left\{\begin{array}{l}

\frac{1}{\omega_{n}(n-2)} \frac{1}{|x-y|^{n-2}}, n \geqslant 3, \\

\frac{1}{2 \pi} \ln \frac{1}{|x-y|}, n=2,

\end{array}\right.

\]



\begin{itemize}

  \item главное фундаментальное решение уравнения Лапласа $\Delta u \equiv \sum_{k=1}^{n} \partial^{2} u / \partial x_{k}^{2}=0$ в $\mathbb{R} n$, где

\end{itemize}



\[

|x-y|=\left[\sum_{k-1}^{n}\left(x_{k}-y_{k}\right)^{2}\right]^{1 / 2}

\]



\begin{itemize}

  \item расстояние между точками $x=\left(x_{1}, \ldots, x_{n}\right)$ и $y=$ $=\left(y_{1}, \ldots, y_{n}\right)$ в $\mathbb{R}^{n}, \omega_{n}=2 \pi^{n / 2} / \Gamma(n / 2)$ - площадь единичной сферы в $\mathbb{R} n, \Gamma$ - гамма-функция. Три интеграла, зависящих от $x$ как от параметра:

\end{itemize}



\[

\left.\begin{array}{rl}

Z(x) & =\int_{G} \rho(y) E(x, y) d y, \\

V(x) & =\int_{S} \mu(y) E(x, y) d S(y), \\

W(x) & =\int_{S} v(y) \frac{\partial}{\partial n_{y}} E(x, y) d S(y),

\end{array}\right\}

\]



где $n_{y}-$ направление внешней относительно $G^{+}$нормали к $S$ в точке $y \in S$, наз. соответственно о 6 ъ е мным потенциалом, потенци алом простогослая и потенци аломдво йного с л о я. Функции $\rho(y), \mu(y)$ и $\boldsymbol{v}(y)$ наз. п ло т нос т я м и соответствующих потенциалов; ниже они будут предполагаться абсолютно интегрируемыми соответственно на $G$ или $S$. При $n=3$ (а иногда и при $n \geqslant 3$ ) интегралы (1) наз. н ь ю то но в ы м о б т м ным по т е нци а ло м, нь ю тоновым и по те нци а лам п постого й во й но гослоя,



\includegraphics[max width=\textwidth, center]{2023_07_08_254cf6ef93625ff1f8feg-1}

аламимас с, простого и двойногослоя.



Пусть $\rho$ принадлежит классу $C^{\mathfrak{1}}(G \cup S)$. Тогда объемный потенциал и его производные 1-го порядка непрерывны всюду в $\mathbb{R} n$, причем их можно вычислить посредством дифференцирования под знаком интеграла, то есть $Z \in C^{1}\left(\mathbb{R}^{n}\right)$. Далее,



\[

\lim _{|x| \rightarrow \infty}(Z(x) / E(x, 0))=M, \quad M=\int_{G} \rho(y) d y .

\]



П роизводные 2-го порядка непрерывны всюду вне $S$, но при переходе через поверхность $S$ они претерпевают разрыв, причем в области $G^{+}$удовлетворяется уравнение Пуассона - $\Delta Z=\rho(x), x \in G^{+}$, а в $G^{-}-$уравнение Лапласа $\Delta Z=0, x \in G^{-}$. Перечисленные свойства характеризуют объемный потенциал.





\end{document}